\documentclass[12pt]{article}

\usepackage[T1]{fontenc}
\usepackage{lmodern}
\usepackage{microtype}
\usepackage[left=1in,right=1in,top=0.82in,bottom=0.84in]{geometry}
\usepackage{setspace}
\usepackage{amsmath,amssymb,mathtools}
\usepackage{mathrsfs}
\usepackage{graphicx}
\usepackage{fancyhdr}
\usepackage[hidelinks]{hyperref}

\hypersetup{
  pdftitle={How the Proof Works},
  pdfauthor={Thomas Birnie}
}

\setstretch{1.23}
\setlength{\parindent}{0pt}
\setlength{\parskip}{0.72\baselineskip}
\setlength{\abovedisplayskip}{15pt plus 3pt minus 2pt}
\setlength{\belowdisplayskip}{15pt plus 3pt minus 2pt}
\setlength{\abovedisplayshortskip}{13pt plus 2pt minus 2pt}
\setlength{\belowdisplayshortskip}{13pt plus 2pt minus 2pt}
\setlength{\jot}{9pt}
\setlength{\emergencystretch}{3em}
\allowdisplaybreaks[2]
\raggedright
\raggedbottom
\clubpenalty=10000
\widowpenalty=10000
\displaywidowpenalty=10000

\pagestyle{fancy}
\fancyhf{}
\renewcommand{\headrulewidth}{0pt}
\fancyfoot[C]{\footnotesize\thepage}

\newcommand{\R}{\mathbb R}
\newcommand{\I}{\mathcal I}
\newcommand{\cR}{\mathcal R}
\newcommand{\Szero}{\mathscr S_0}

\begin{document}

{\Large\bfseries How the Proof Works\par}
\vspace{0.16\baselineskip}
{\small Thomas Birnie\par}

\vspace{0.55\baselineskip}

Finite kinetic energy does not rule out concentration. An incompressible
vortex can stretch along its axis only by narrowing across it. If \(r\) is
the core radius and \(U\) is the change in velocity across that core, then
\[
  |\nabla u|
  \sim
  \frac{U}{r},
  \qquad
  \int_{B_r}|u|^2\,dx
  \sim
  U^2r^3.
\]
Shrinking \(r\) raises the velocity gradient through the factor \(1/r\), while
the energy carried by that region can remain finite because its volume falls
like \(r^3\). The energy bound leaves that possibility open, so the equation
has to tell us what happens as the scale shrinks:

\[
  \partial_tu
  +
  (u\cdot\nabla)u
  +
  \nabla p
  =
  \nu\Delta u,
  \qquad
  \nabla\cdot u=0.
\]
Transport can sharpen the profile, while viscosity diffuses it. Rescaling a
concentrating profile shows how both effects change as the profile narrows. For
\(\lambda>1\), set
\[
  u_\lambda(x,t)
  =
  \lambda u(\lambda x,\lambda^2t),
  \qquad
  p_\lambda(x,t)
  =
  \lambda^2p(\lambda x,\lambda^2t).
\]
The spatial width is divided by \(\lambda\), and the time scale is divided by
\(\lambda^2\). With
\(\Phi_\lambda(x,t)=(\lambda x,\lambda^2t)\), the chain rule gives
\[
\begin{aligned}
  \partial_tu_\lambda
  &=
  \lambda^3(\partial_tu)\circ\Phi_\lambda,
  &
  (u_\lambda\cdot\nabla)u_\lambda
  &=
  \lambda^3((u\cdot\nabla)u)\circ\Phi_\lambda,
  \\
  \nabla p_\lambda
  &=
  \lambda^3(\nabla p)\circ\Phi_\lambda,
  &
  \nu\Delta u_\lambda
  &=
  \lambda^3(\nu\Delta u)\circ\Phi_\lambda.
\end{aligned}
\]
All four parts grow by \(\lambda^3\). A narrower profile strengthens nonlinear
transport and viscous diffusion together; scaling alone does not give either
one control of the small scales.

Under this rescaling,
\[
\begin{aligned}
  \widehat{u_\lambda}(\xi,t)
  &=
  \lambda^{-2}
  \widehat u(\xi/\lambda,\lambda^2t),
  \\
  \|u_\lambda(t)\|_{\dot H^s}^2
  &=
  \int_{\R^3}
  |\xi|^{2s}\lambda^{-4}
  |\widehat u(\xi/\lambda,\lambda^2t)|^2\,d\xi
  \\
  &=
  \lambda^{2s-1}
  \|u(\lambda^2t)\|_{\dot H^s}^2.
\end{aligned}
\]
For 
\(s<\tfrac12\), the norm falls as the profile narrows. For
\(s>\tfrac12\), it rises. At
\(s=\tfrac12\), it is unchanged. So \(s=\tfrac12\) is critical: its norm
neither grows nor shrinks as the profile narrows.

At that height, the critical measurement stays unchanged when the rescaling is
repeated. Set \(\lambda=2\). Halving the core radius then quarters the
corresponding time scale:
\[
  r_n=2^{-n}r_0,
  \qquad
  \delta t_n=4^{-n}\delta t_0,
  \qquad
  \sum_{n=0}^{\infty}\delta t_n
  =
  \frac43\delta t_0.
\]
The series shows how infinitely many scale reductions could fit inside a finite
time. A singularity would require one solution to pass through every scale in
that sequence.

In order for a singularity to form, velocity has to blow up. In order for
that to happen, its acceleration would also have to blow up, which means its
jerk has to blow up, and so on, and so on:
\[
  u,
  \qquad
  \partial_tu,
  \qquad
  \partial_t^2u,
  \qquad
  \ldots
\]
If a singularity were forming, each higher time derivative should show less
spatial regularity, not more. The differentiated equations reverse that
expectation.

\begingroup
\setlength{\abovedisplayskip}{17pt}
\setlength{\belowdisplayskip}{17pt}
\begin{flushleft}
\resizebox{\textwidth}{!}{%
\(
\begin{aligned}
  \partial_tu
  &=
  \nu\Delta u-(u\cdot\nabla)u-\nabla p,
  \\[0.95\baselineskip]
  \partial_t^2u
  &=
  \nu\Delta\partial_tu-(\partial_tu\cdot\nabla)u-(u\cdot\nabla)\partial_tu-\nabla\partial_tp,
  \\[0.95\baselineskip]
  \partial_t^3u
  &=
  \nu\Delta\partial_t^2u-(\partial_t^2u\cdot\nabla)u-2(\partial_tu\cdot\nabla)\partial_tu-(u\cdot\nabla)\partial_t^2u-\nabla\partial_t^2p.
\end{aligned}
\)
}
\end{flushleft}
\endgroup

For \(n\geq0\),
\[
\begin{aligned}
  \partial_t^{\,n+1}u
  ={}&
  \nu\Delta\partial_t^nu
  \\
  &-
  \sum_{a=0}^{n}
  \binom{n}{a}
  \bigl(\partial_t^au\cdot\nabla\bigr)
  \partial_t^{\,n-a}u
  \\
  &-
  \nabla\partial_t^np.
\end{aligned}
\]
Set
\[
  V_n:=\partial_t^nu,
  \qquad
  Q_n:=\partial_t^np.
\]

At the critical height, that temporal sequence appears inside one energy. Write
\[
  H(t)
  :=
  \frac12\|\Lambda^{1/2}u(t)\|_2^2,
  \qquad
  \Lambda:=(-\Delta)^{1/2}.
\]
For any \(s\geq0\), set
\[
  E_s(t)
  :=
  \frac12\|\Lambda^su(t)\|_2^2.
\]
Let \(\mathcal P_s\) denote the transport-pressure contribution at height
\(s\). Pairing the viscous term at that height gives
\[
  \nu\operatorname{Re}
  \left\langle
    \Lambda^s\Delta u,
    \Lambda^su
  \right\rangle_{L^2}
  =
  -\nu\|\Lambda^{s+1}u\|_2^2
  =
  -2\nu E_{s+1}.
\]
Adding \(\mathcal P_s\) gives
\[
  E_s'
  =
  \mathcal P_s
  -
  2\nu E_{s+1}.
\]
This is where viscosity moves the calculation from height \(s\) to height
\(s+1\), with coefficient \(2\nu\).

With \(H=E_{1/2}\),
\[
\begin{aligned}
  H'
  &=
  \mathcal P_{1/2}-2\nu E_{3/2},
  \\
  H''
  &=
  \partial_t\mathcal P_{1/2}
  -2\nu\mathcal P_{3/2}
  +(2\nu)^2E_{5/2},
  \\
  H'''
  &=
  \partial_t^2\mathcal P_{1/2}
  -2\nu\partial_t\mathcal P_{3/2}
  +(2\nu)^2\mathcal P_{5/2}
  -(2\nu)^3E_{7/2}.
\end{aligned}
\]
At order \(n\geq1\),
\[
  H^{(n)}
  =
  (-2\nu)^nE_{n+1/2}
  +
  \sum_{j=0}^{n-1}
  (-2\nu)^j
  \partial_t^{\,n-1-j}\mathcal P_{j+1/2}.
\]
The \(n\)-th derivative of \(H\) reaches \(E_{n+1/2}\). As temporal order
rises, Sobolev height rises with it:
\[
  \text{higher temporal order}
  \quad\longrightarrow\quad
  \text{higher spatial Sobolev order}.
\]
Control of every finite height through \(T_*\) would give, by Sobolev
embedding,
\[
  s>k+\frac32
  \quad\Longrightarrow\quad
  H^s(\R^3)
  \hookrightarrow
  C^k(\R^3),
\]
and hence
\[
  \bigcap_{m<\infty}H^m(\R^3)
  =
  H^\infty(\R^3)
  \hookrightarrow
  C^\infty(\R^3).
\]

Fix finite temporal order \(N\) and finite Sobolev height \(M\). At temporal
order \(r\),
\[
\begin{aligned}
  E_{r,s}(t)
  &:=
  \frac12\|\Lambda^sV_r(t)\|_2^2,
  \\
  \partial_tE_{r,s}
  &=
  \mathcal P_{r,s}
  -
  2\nu E_{r,s+1}.
\end{aligned}
\]
Stopping at \(r=N\) and \(s=M\) leaves finitely many adjacent rows. Their total
velocity norm \(\mathfrak R_{N,M}(u;T)\) stays bounded as \(T\uparrow T_*\):
\[
  \sup_{0<T<T_*}
  \mathfrak R_{N,M}(u;T)
  <
  \infty
  \qquad (N,M<\infty).
\]
Adding the corresponding \(Q_0,\ldots,Q_N\) coordinates gives the
pressure-complete rectangle \(\cR_{N,M}[u_0;T]\).

Whenever two rectangles overlap, their common entries are derivatives of \(u\)
and \(p\) already fixed by the differentiated equations, so the overlaps
agree. Their compatible intersection is
\[
  \I[u_0]
  :=
  \bigl(\cR_{N,M}[u_0]\bigr)_{N,M<\infty},
\]
and every finite pair of indices occurs in one of those blocks. Hence
\[
  u
  \in
  \bigcap_{r,m<\infty}
  H^r(0,T_*;H^m(\R^3)).
\]

Projecting onto \(r=0\) gives
\[
  \Szero
  :=
  \pi_0\I[u_0].
\]
For every finite spatial height,
\[
  u\in C([0,T_*];H^m(\R^3)).
\]
Thus
\[
  u_*:=u(T_*)
  \in
  \bigcap_{m<\infty}H^m(\R^3)
  =
  H^\infty(\R^3)
  \hookrightarrow
  C^\infty(\R^3).
\]
Local existence gives a smooth solution starting from \(u_*\), and uniqueness
makes it agree with \(u\) on their overlap. The solution passes through
\(T_*\), contradicting finite maximality. Thus
\[
  T_*=\infty.
\]

For a finite \(L^p\) norm of \(\nabla^k\partial_t^ru\), choose an integer
\(m\) with
\[
  m>k+\frac32-\frac3p,
  \qquad
  2\leq p\leq\infty.
\]
Since
\(\partial_t^ru\in C_tH_x^m\), Sobolev embedding gives
\[
  \nabla^k\partial_t^ru
  \in
  L_t^\infty L_x^p.
\]
Choosing \(m=1,2,3\) gives
\[
\begin{aligned}
  \I[u_0]
  &\longrightarrow
  C_tH_x^1
  \longrightarrow
  L_t^\infty L_x^6,
  \\
  \I[u_0]
  &\longrightarrow
  C_tH_x^2
  \longrightarrow
  L_t^\infty L_x^\infty,
  \\
  \I[u_0]
  &\longrightarrow
  C_tH_x^3
  \longrightarrow
  L_t^\infty W_x^{1,\infty}.
\end{aligned}
\]
For \(m=3\),
\[
\begin{aligned}
  \int_0^{T_*}\|\nabla u(t)\|_\infty\,dt
  &\leq
  T_*
  \sup_{0\leq t\leq T_*}
  \|\nabla u(t)\|_\infty
  \\
  &\leq
  C T_*
  \sup_{0\leq t\leq T_*}
  \|u(t)\|_{H^3}
  <
  \infty.
\end{aligned}
\]

In shorthand,
\[
\begin{aligned}
  u_0
  &\longrightarrow
  \{(V_n,Q_n)\}_{n\geq0}
  \longrightarrow
  \{\cR_{N,M}[u_0]\}_{N,M<\infty}
  \\
  &\longrightarrow
  \I[u_0]
  \longrightarrow
  \Szero
  \longrightarrow
  u_*\in H^\infty
  \longrightarrow
  T_*=\infty.
\end{aligned}
\]
The estimates come from the intersection, not the other way around.

\end{document}
